% --- File: sections/02_regulatory.tex ---

The deployment of machine learning in institutional risk management is strictly governed by international capital adequacy standards, most notably the Basel IV framework and the Fundamental Review of the Trading Book (FRTB). To understand the necessity of an ``Anchored'' neural architecture, one must first examine the punitive nature of regulatory backtesting.

Under the FRTB framework \cite{bcbs2019frtb}, the primary capital calculation has migrated from a 99\% Value-at-Risk (VaR) measure to a 97.5\% Expected Shortfall (ES, also known as CVaR). ES is theoretically superior: as a coherent risk measure, it captures the \emph{severity} of losses beyond the threshold, not merely the probability of exceeding it. A model producing rare but catastrophic losses can satisfy a VaR constraint while exhibiting unacceptable ES---a structural gap that regulators have explicitly moved to close.

\textbf{Why, then, does this paper focus on VaR?} The answer lies in a fundamental statistical constraint. ES is an expectation over the tail of the distribution and cannot be directly validated through classical backtesting. To apply the Kupiec Proportion of Failures test---or any breach-count framework---one needs a binary daily event: either the realized return exceeded the forecast threshold or it did not. VaR provides exactly this. ES does not: there is no clean definition of a single-day ``ES breach.'' Accordingly, the FRTB framework retains daily 99\% VaR backtesting as the \emph{regulatory approval gate}. An institution must first demonstrate that its VaR model is well-calibrated before it is permitted to use the IMA and apply ES-based capital charges. Failing the Kupiec test on VaR disqualifies the model entirely, regardless of its ES properties. In this sense, VaR is not obsolete---it is the entrance exam.

From a methodological standpoint, solving the VaR problem first is also the correct architectural sequence. The Asymmetric Quantile Loss (Pinball Loss) at quantile $q = 0.01$ directly targets the 99\% VaR. Expected Shortfall is the integral of quantiles above that threshold: $\text{ES}^q = \mathbb{E}[r \mid r < \text{VaR}^q]$. The quantile regression framework developed here is therefore the natural foundation for a full ES architecture, and extending it to a 97.5\% ES target via a Mixture Density Network is identified as the primary direction for future work (Section 6).

Regulators evaluate model validity over a rolling 250-day window using a zone-based framework derived from the Kupiec Proportion of Failures (POF) test \cite{kupiec1995techniques}: too few breaches signal over-conservatism that traps capital, while too many signal model failure and can ultimately revoke the institution's Internal Models Approach (IMA) approval in favour of the more punitive Standardized Approach. In this study, we adopt the Kupiec POF test directly as our acceptance criterion, targeting a breach rate statistically consistent with the 1\% theoretical exceedance probability.

Under conditions of macroeconomic stress and for highly volatile, fat-tailed asset classes, this regulatory structure creates a critical optimization problem. Traditional models such as Historical Simulation often react too slowly to regime shifts, accumulating excess breaches during sudden market downturns. Conversely, to stay safely within the Kupiec acceptance region, risk managers frequently over-calibrate parametric models --- reserving artificially large capital buffers that drag down Return on Equity (ROE).

Pure Deep Learning models often exacerbate this issue rather than solving it. A naive LSTM, lacking structural constraints, can suffer from ``catastrophic forgetting'' or overfit to low-volatility training data. When deployed out-of-sample, a sudden volatility spike can rapidly accumulate exceptions, causing the model to fail the Kupiec test and jeopardize the institution's IMA status.

Therefore, a practically deployable neural architecture must not only capture non-linear market dynamics but must do so within guaranteed statistical bounds. The objective is to pass the Kupiec test at the correct calibration level --- neither too conservative nor too permissive --- while minimizing the mean VaR to maximize capital efficiency. This regulatory necessity motivates the asymmetric Quantile Loss function and parametric anchoring detailed in Section 3.